\section{Introduction}\label{sec:introduction}

A matching in a graph is a set of edges with pairwise disjoint endpoints.
The maximum-cardinality matching problem asks for such a set with the largest
possible number of edges. We prove an almost-linear bound in the adjacency-list
input size for general graphs, with explicit output and a single uniform algorithm.
The running-time guarantee includes every unsuccessful random computation,
all finite-precision work, and all input and output costs.

\subsection{Statement and computational model}

The input is a simple undirected unweighted graph $G=([n],E)$, where $n\ge1$
and $m=|E|$. It is given by an array of adjacency-list boundaries and two
neighbor occurrences per edge. Labels and indices use $O(\log(n+2))$ bits.
The lists may have any valid order. Isolated vertices count toward $n$, and
the graph need not be connected. Write $p=n+m$, which is within a constant
factor of the number of input word records, and write $\nu(G)$ for the maximum
matching size.

We use a classical randomized word-RAM whose word length is
\[
                     w=\lceil c_0\log_2(n+2)\rceil,
\]
for a fixed constant $c_0\ge2$ chosen with the algorithm. Unit-cost instructions
are word reads and writes, comparison, addition, subtraction, multiplication
modulo $2^w$, integer quotient and remainder with nonzero divisor, Boolean
operations, and shifts by fewer than $w$ positions. Each independent fair
random bit costs one instruction. Longer arithmetic uses charged multiword
implementations. Allocation, initialization, data movement, preprocessing,
and output writing are included. There is no oracle or nonuniform advice.
Since $n\le p\le n^2$ for a simple graph, the word length is also
$\Theta(\log(p+2))$.

\begin{theorem}[Exact matching in almost-linear input-size time]\label{thm:main}
There are a single uniform randomized word-RAM algorithm, fixed constants
$c_0\ge2$ and $C\ge1$, and a nonincreasing function
$\eta:\mathbb N_{\ge1}\to\mathbb R_{\ge0}$ with $\eta(q)\to0$ as
$q\to\infty$ such that the following holds.
For every input graph $G=([n],E)$ with $m=|E|$ and $p=n+m$, the algorithm
halts on every computation path in at most
\[
                          C p^{1+\eta(p)}
\]
instructions. With probability at least $2/3$ it outputs an explicit list of
edges of a matching of size $\nu(G)$.
\end{theorem}

For every fixed integer $\kappa\ge1$, this gives
$C_\kappa n^{1+\eta(n)}$ instructions on graphs with $m\le\kappa n$, for
some constant $C_\kappa$. The same algorithm and exponent envelope work for
all $\kappa$, which need not be supplied as an input. At the dense extreme
$m=\Theta(n^2)$, the bound is $n^{2+o(1)}$.

A perfect matching is not promised. In particular, the perfect-matching
procedure for graphs of maximum degree three will be used only after
reductions that preserve explicit witnesses. A separate algorithm for each fixed exponent slack
would also be insufficient; Section~\ref{sec:uniformization} constructs the
single finite program and the common envelope in Theorem~\ref{thm:main}.

A reduction also gives an almost-linear bound for prescribed vertex degrees. Given an
integer $f(v)$ with $0\le f(v)\le\deg_G(v)$ for every vertex $v$, an
\emph{$f$-factor} is a spanning subgraph $([n],F)$ with
$\deg_F(v)=f(v)$ for every $v$. Corollary~\ref{cor:f-factor} gives one
uniform randomized algorithm that decides whether an $f$-factor exists and
returns its edge list when it does, with success probability at least $2/3$
and $(n+m)^{1+o(1)}$ instructions on every computation path. The degree
prescription is an array of one-word integers. Section~\ref{sec:f-factors}
gives the witness-preserving reduction with only logarithmic expansion in
the number of edge incidences.

\subsection{Context and significance}

The structure of matchings has long connected combinatorial existence
criteria with efficient construction. Hall's theorem characterizes when
a family of sets has distinct representatives~\cite{Hall1935}, and Tutte's
perfect-matching theorem identifies the obstruction caused by odd
components~\cite{Tutte1947}. Berge's augmenting-path criterion turns
maximum cardinality into a question about alternating
paths~\cite{Berge1957}. Edmonds made this approach algorithmic in general
graphs through blossom contraction and expansion~\cite{EdmondsPaths}.
His matching-polytope theorem gives a complementary explanation of the
odd-set constraints~\cite{EdmondsPolytope}: vertex-degree inequalities
alone allow, for example, value $1/2$ on each edge of a triangle.

Tutte's short proof of the factor theorem reduces prescribed-degree
$f$-factors to perfect matchings~\cite{Tutte1954Factor}. The reduction in
Section~\ref{sec:f-factors} follows this incidence-port approach, replacing
the local complete connections to auxiliary vertices by the rearrangeable
networks used in Section~\ref{sec:uniformization}. This controls the size of
the perfect-matching instance even when a host vertex has high degree.
There are also direct algorithms: Gabow finds a maximum-cardinality edge
set subject to $\deg_F(v)\le f(v)$ in $O(n^{2/3}m)$ time on simple
graphs~\cite[Theorem~4.1]{Gabow2025FMatching}. After linear-time input
processing, including removal of zero-demand vertices, this gives an
$O(n+n^{2/3}m)$ factor test: an $f$-factor exists exactly when the returned
set has $2|F|=\sum_v f(v)$.

Batching shortest augmenting paths led to the bipartite algorithms of
Hopcroft and Karp~\cite{HopcroftKarp1973} and, independently,
Karzanov~\cite{Karzanov1973Representatives}, and to the Micali--Vazirani
algorithm for general graphs~\cite{MicaliVazirani,MicaliVaziraniAnalysis}.
The latter has running time $O(m\sqrt n)$ on a RAM; reading the explicit
input and removing isolates gives the convenient universal bound
$O(n+m\sqrt n)$. This is $O_\kappa(n^{3/2})$ when $m\le\kappa n$.
Gabow and Tarjan developed a different route through weighted
scaling~\cite{GabowTarjan1991}, and Gabow later gave a direct
weighted-search formulation at the same cardinality bound~\cite{Gabow2017}.
On very dense graphs, Goldberg and Karzanov improve the combinatorial bound to
$O(n+m\sqrt n\log(n^2/m)/\log n)$ by combining skew-symmetric flow
with graph compression~\cite[Theorem~10.3]{GoldbergKarzanov2004}.
The weighted problem has its own sequence of improvements: the scaling
algorithm of Duan, Pettie, and Su finds an exact maximum-weight matching in
$O(m\sqrt n\log(nW))$ time for integer weights of magnitude at most
$W$~\cite{DuanPettieSu2018}, with logarithms bounded below by a constant.

Algebraic methods offer another important approach. Random evaluation
of the Tutte matrix connects matching to linear algebra~\cite{Lovasz1979}.
Mucha and Sankowski, and subsequently Harvey, obtained randomized
algorithms that construct maximum matchings at the matrix-multiplication
time scale, conventionally written $O(n^\omega)$
field operations~\cite{MuchaSankowski2004,Harvey2009}.
These algorithms already recover explicit witnesses and can work over
finite fields of polynomial size. Their matrix-multiplication bound is
most competitive on dense graphs. Theorem~\ref{thm:main} instead gives an
almost-linear bound in the explicit input size throughout the range of
densities, including $n^{2+o(1)}$ at the dense extreme.

For bipartite matching, the unit-capacity flow formulation permits
continuous optimization to accelerate exact search. M\k{a}dry's
electrical-flow interior-point method~\cite{Madry2013CentralPath} and
the moderately dense algorithm of van den Brand
et al.~\cite{vanDenBrandEtAl2020BipartiteMatching} were intermediate
advances. The almost-linear exact integral-flow algorithms of Chen
et al.~\cite{FlowRandomized} and van den Brand
et al.~\cite{FlowDeterministic} then yield almost-linear exact bipartite
matching. General graphs require additional control of odd blocks.
Here that control is provided by the matching state and its local
projections. The deterministic flow theorem is used through the
explicit static interface in Lemma~\ref{lem:integral-flow}.

The proof combines several established ideas with interfaces tailored
to that state. Its constructive factor-critical blocks belong to the
contraction-and-expansion tradition, while Hall's theorem underlies
the reassignment of their exposed vertices. The multiscale allocation
and debt bounds are proved directly. The routing helper uses the
cut-matching structure of Khandekar, Rao, and
Vazirani~\cite{KhandekarRaoVazirani2006}, together with concentration
and conductance estimates~\cite{TroppBernstein,SinclairJerrum1989}.
Its recursive virtual graphs follow the hierarchical random-walk
embeddings of Ghaffari, Kuhn, and Su~\cite[Section~3.1.2]{GhaffariKuhnSu2017}.
Here the sampled support and the substituted routes must also preserve
supplied resource loads, with their full construction costs charged.

The use of approximate minimum-ratio circulations as optimization
directions has its modern antecedent in the almost-linear flow
framework~\cite{FlowRandomized,FlowDeterministic}. Chen et al.
prove a deterministic \emph{fully dynamic} approximate
minimum-ratio-cycle theorem with flat forest
representations~\cite[Definition~2.3 and Theorem~3.5]{IncrementalGraphs}.
Our routing hierarchy is related to the origin-aware construction of
Rozho\v{n} et al.~\cite{MinorAggregateRouting}; the finite forest family
and associated cycle extraction adapt the construction in
\cite[Section~4.5]{IncrementalGraphs}. The preallocated forest copies used
to lift recursive paths follow Kyng, Meierhans, and Probst
Gutenberg~\cite[Section~3.3, Lemma~3.21 and Algorithm~5]{KyngMeierhansProbstGutenberg2023FlatForest}.
Sections~\ref{sec:engine} and~\ref{sec:backend-updates} specify how retained
mapped edges preserve accumulated coordinate changes and absolute motion
through these representation updates. The recent distance-oracle
approach of Kyng et al.~\cite{MinRatioDistanceOracles2026} provides
another development of this method. Sections~\ref{sec:backend-quality}
and~\ref{sec:backend-updates} supply the complete backend used here,
including its adaptive guarantees, explicit changes of representation,
stable edge-copy lifetimes, and bounds on traversal length before
cancellation. Those guarantees are the inputs to our priced-step
analysis.

\subsection{Main ideas and organization}

The central subproblem is to find a perfect matching in a graph of maximum
degree three.  Section~\ref{sec:uniformization} reduces exact cardinality
queries to this subproblem while retaining the original edge identifiers,
so that every successful answer can be decoded into a matching of the
input graph.

The perfect-matching procedure maintains a current matching through a
partition of the vertices.  Its odd blocks, called \emph{bins}, come with
internal matchings after any prescribed vertex is left exposed.
Some outside singleton vertices, called \emph{rows}, are assigned along
original edges to distinct bins.  Each assigned bin matches its exposed
vertex to its row; an unassigned bin contributes one unmatched vertex.
The remaining vertices are stored as matched pairs, apart from unmatched
vertices set aside during repairs.  This representation permits a change
of exposure without immediately expanding a nested block.

The difficulty is to repair this representation after structural changes
without repeatedly scanning all its odd blocks. A \emph{ledger} allocates
a fixed amount from each row among its neighboring bins. The difference
between a bin's capacity and its allocated load is its residual. Several
ledgers impose residual margins proportional to bin size; a nonnegative
\emph{debt} allows a margin to fail temporarily. At each settled state,
the lowest ledger has no debt and certifies a strict Hall surplus for
every nonempty set of bins. The highest ledger carries a logarithmic
objective that must admit progress when a perfect matching exists and no
structural operation is available.

The intermediate ledgers turn that progress into affordable repairs. A
repair is an exact flow problem on row--bin incidences, initially explored
only at bins whose debt is substantial relative to their size. An
unexplored boundary bin temporarily accepts an optimistic amount of flow.
If enough arrives, the algorithm explores that bin and permanently locks
in an absorbed amount proportional to its size. The total available flow
is bounded by the starting debt, so these locks bound the number of
physical vertices examined. If the requested residual margin is
infeasible, the repair instead removes a set of bins. The gap between
adjacent margins makes this removal decrease debt in every preserved
coarser ledger by an amount proportional to the removed size. Thus debt
pays both for local exploration and for removing blocks from the state.
The logarithmic optimization uses exact
conservation throughout; integral cut inequalities turn a point whose
total bound violation is below one into a certificate of exact
feasibility.  A final local cut also recovers the entire residual
source-reachable set needed by the matching repair.

Both the matching objective and these local flow tests use the same
\emph{priced-step} interface.  It searches for a circulation direction
with sufficient objective increase relative to its weighted length.
When it accepts such a direction, it takes a small exact dyadic step.  The direction is represented by
individual edges and paths in a changing forest, so a long path need not
be expanded for every step.  The implementation separately records
signed coordinate changes and absolute motion.  The latter counts
repeated traversals before cancellation and controls when a coordinate
must be read and repriced.

The cycle data structure supplying this interface is constructed in
full.  A static routing sparsifier first provides short routes whose
total use of specified resources is bounded.  During an interval of
updates, a frozen collection of forests then decomposes every current
circulation into local correction cycles and a circulation in one
smaller graph.  A hierarchy of routing snapshots limits the number of
new representative edges added to that graph.  Finally, prebuilt copies
of the paths represented by its edges lift its forest back to the
current graph.  When recursive occurrences change at otherwise unaffected
child labels, these copies retain their identities. This preserves the exact
accumulated values used by the priced-step implementation; replacing an affected copy first transfers its saved history.

Let $S$ be a power-of-two upper bound on the number of vertices in the
reduced graph. For each sufficiently large fixed parameter $a$, these
ingredients yield time $S^{1+O(1/a)+o_a(1)}$ for a subcubic
perfect-matching test. The coefficient
hidden in $O(1/a)$ is absolute. Logarithmic powers may depend on the fixed
parameter. Deterministic prerequisite failures suspend a parameter branch;
exhausted random retries terminate its test with a possibly erroneous
negative answer. This distinction permits an every-path asymptotic bound.
Finally, instruction-level dovetailing with summable error budgets turns all
fixed-parameter guarantees into Theorem~\ref{thm:main}, without a convergence
rate or an input-dependent table of thresholds.

Sections~\ref{sec:matching-state} and~\ref{sec:projection-analysis} develop
the matching algorithm assuming the priced-step interface.
Section~\ref{sec:engine} proves that interface from the flat-cycle backend.
Section~\ref{sec:routing-helper} gives the static helper, and
Sections~\ref{sec:backend-quality} and~\ref{sec:backend-updates} prove
the backend's quality and implementation guarantees.
Section~\ref{sec:uniformization} completes the reductions and the uniform
finite-word implementation.
Section~\ref{sec:f-factors} applies the resulting matching algorithm to
prescribed-degree factors through a linearithmic-size reduction.

The principal interfaces and their proof dependencies are summarized below.
\begin{center}
\begin{tabular}{@{}p{.22\textwidth}p{.73\textwidth}@{}}
\toprule
Result & Role in the proof \\
\midrule
Theorem~\ref{thm:subcubic} & Fixed-parameter subcubic perfect-matching search;
  Sections~\ref{sec:matching-state}--\ref{sec:projection-analysis} use
  Lemma~\ref{lem:priced-steps}. \\[3pt]
Lemma~\ref{lem:priced-steps} & Priced logarithmic steps;
  Section~\ref{sec:engine} proves the interface from
  Lemma~\ref{lem:flat-cycle}. \\[3pt]
Lemma~\ref{lem:flat-cycle} & Flat-cycle backend;
  Sections~\ref{sec:backend-quality}--\ref{sec:backend-updates} construct it
  using the static routing helper, Lemma~\ref{lem:routing-helper}. \\[3pt]
Theorem~\ref{thm:main} & Explicit maximum matching with one input-size exponent
  envelope; Section~\ref{sec:uniformization} applies the witness reductions
  and instruction-level schedule to Theorem~\ref{thm:subcubic}. \\
\bottomrule
\end{tabular}
\end{center}
